\chapter{取引コストと執行}
\label{ch:cost}

論文の戦略の多くは、取引コストを引く前の成績で示される。実際の運用では、コストが予測の価値の大半を食うことが珍しくない。クオンツの運用会社の競争力の大きな部分は、予測の質よりも、コストを測り、減らす技術にある。

\section{コストの内訳}

\begin{table}[H]
\centering\small
\caption{取引コストの内訳}
\label{tab:costs}
\begin{tabular}{@{}P{0.24\linewidth}P{0.46\linewidth}P{0.22\linewidth}@{}}
\toprule
コスト & 中身 & 見え方 \\
\midrule
手数料・税 & 証券会社への手数料、取引所の手数料、取引税 & 明示的 \\
スプレッド & 買値と売値の差の半分（すぐに売買すると払う） & 半ば明示的 \\
マーケット・インパクト & 自分の売買が価格を不利に動かす分 & 暗黙 \\
遅延 & 判断してから売買するまでの価格の変化 & 暗黙 \\
機会の費用 & 売買しきれなかった分の、逃した利益 & 暗黙 \\
空売りの費用 & 借りた株の貸株料、配当の支払い、リコールのリスク & 明示的 \\
資金の費用 & 借入れの金利、証拠金に寝る資金 & 明示的 \\
\bottomrule
\end{tabular}
\end{table}

大型株を少し売買するなら、コストの大部分はスプレッドで、片道 0.01〜0.05\% 程度である。小型株や大口の売買では、マーケット・インパクトが支配的になり、片道 0.5\% を超えることも珍しくない。

\section{インプリメンテーション・ショートフォール}

Perold（1988）\cite{perold1988} は、コストを「売買を決めた時点の価格で、すべて売買できたとした紙の上のポートフォリオ」と「実際のポートフォリオ」の成績の差として測ることを提案した（インプリメンテーション・ショートフォール）。
\begin{equation}
\mathrm{IS} = \underbrace{\sum_j q_j\,(p_j - p_0)}_{\text{実際に売買した分のコスト}} + \underbrace{(Q - \textstyle\sum_j q_j)(p_T - p_0)}_{\text{売買しなかった分の機会の費用}} + \text{手数料}
\end{equation}
（買いの場合。$p_0$ は判断の時点の価格、$q_j$, $p_j$ は各約定の数量と価格、$Q$ は目標の数量）。この物差しは、手数料・スプレッド・インパクト・遅延・機会の費用をすべて含む。

\section{マーケット・インパクト}

\subsection{一時的なインパクトと恒久的なインパクト}

大口の注文を出すと、価格は不利に動く。
\begin{itemize}
\item \textbf{一時的なインパクト}：板の流動性を消費したことによる一時的な価格の動き。売買を終えると、ある程度戻る。
\item \textbf{恒久的なインパクト}：売買が情報を含むと市場に解釈され、価格の水準が変わる分。戻らない。
\end{itemize}

Kyle（1985）\cite{kyle1985} のモデルでは、情報を持つ投資家の注文の流れに対して、価格は線形に動く（$\Delta p = \lambda\,Q$、$\lambda$ は「Kyle のラムダ」）。

\subsection{平方根則}

実際の大口の注文（何日かに分けて執行する「メタオーダー」）のデータでは、インパクトは注文の大きさの平方根にほぼ比例することが、多くの市場・時代で観察されている \cite{athl2005,toth2011,bouchaud2018}。
\begin{equation}
I(Q) \approx Y\,\sigma\,\sqrt{\frac{Q}{V}}
\label{eq:sqrt}
\end{equation}
$\sigma$ は日次のボラティリティ、$V$ は 1 日の売買高、$Q$ は注文の量、$Y$ は 1 前後の定数。インパクトは、取引額そのものではなく、\textbf{売買高に対する割合}と、\textbf{ボラティリティ}で決まる。

\begin{exbox}[平方根則で見積もる]
\begin{itemize}
\item 大型株：日次のボラティリティ 1.5\%、1 日の売買高 5 億ドル。100 万ドルの買い（売買高の 0.2\%）：$1.5\% \times \sqrt{0.002} \approx 0.07\%$。
\item 小型株：日次のボラティリティ 3\%、1 日の売買高 500 万ドル。同じ 100 万ドルの買い（売買高の 20\%）：$3\% \times \sqrt{0.2} \approx 1.3\%$。
\end{itemize}
同じ金額でも、小型株ではインパクトが 20 倍近くになる。小型株の戦略が「論文では儲かるのに、大きな資金では儲からない」理由の一つである。
\end{exbox}

\section{最適執行}

\subsection{Almgren--Chriss のモデル}

大口の注文を、どのくらいの速さで執行するか。速く売買するとインパクトが大きく、ゆっくり売買すると、その間の価格の変動のリスクを負う。Almgren・Chriss（2000）\cite{almgrenchriss2000} は、このトレードオフを平均分散の問題にした。

時間 $[0, T]$ で $X$ 株を売る。残りの株数を $x(t)$、売る速さを $v = -\dot x$ とし、一時的なインパクトが速さに比例する（$\eta v$）とすると、コストの期待値とリスクは
\begin{equation}
\E[\text{コスト}] = \int_0^T \eta\,v(t)^2\,dt + (\text{恒久的なインパクトの項}), \qquad
\Var[\text{コスト}] = \sigma^2\int_0^T x(t)^2\,dt
\end{equation}
で、$\E + \lambda\Var$ を最小にする執行の軌跡は
\begin{equation}
x(t) = X\,\frac{\sinh\big(\kappa(T-t)\big)}{\sinh(\kappa T)}, \qquad \kappa = \sqrt{\frac{\lambda\sigma^2}{\eta}}
\end{equation}
となる。リスクを嫌うほど（$\lambda$ が大きいほど）、最初に速く売る。$\lambda \to 0$ なら一定の速さ（TWAP）になる。

\subsection{執行のアルゴリズム}

実務の執行アルゴリズムは、証券会社や運用会社が提供する。

\begin{itemize}
\item \textbf{VWAP}：その日の売買高の時間ごとの分布に合わせて執行し、その日の売買高加重の平均価格に近づける。
\item \textbf{TWAP}：時間で均等に分ける。
\item \textbf{POV}（参加率）：市場の売買高の一定の割合（例：10\%）で執行する。
\item \textbf{IS}（ショートフォール）：判断の時点の価格からの差を最小にする（Almgren--Chriss の考え方）。
\item \textbf{引けの売買}：取引所の終値を決める板寄せ（クロージング・オークション）で売買する。指数のファンドの売買が集中し、流動性が大きい。
\item \textbf{ダークプール・スマートオーダールーティング}：公開されない板や複数の取引所に注文を振り分け、情報の漏れとインパクトを減らす。
\end{itemize}

\section{コストを入れた最適化}

ポートフォリオの最適化（\chapref{ch:pc}）に、コストを直接入れる。
\begin{equation}
\max_{\bm w}\ \bm w\T\bm\alpha - \frac{\lambda}{2}\bm w\T\Sigma\bm w - \sum_i \Big( c_i\,|w_i - w_i^0| + k_i\,|w_i - w_i^0|^{3/2} \Big)
\end{equation}
第 1 の項は比例するコスト（スプレッド・手数料）、第 2 の項は平方根則のインパクト（取引額 $\times$ 平方根 $=$ 1.5 乗）。比例するコストがあると、予測が少し変わっただけでは売買しない「売買しない領域」ができる。これで回転率が大きく下がり、コストの後の成績はむしろ良くなる。

\section{容量}

戦略の\textbf{容量}（キャパシティ）は、コストの後の利益が最大になる、あるいはゼロになる運用資金の大きさである。資金 $A$ の戦略で、取引額は $A$ に比例し、平方根則のインパクトの率は $\sqrt{A}$ に比例するので、
\begin{equation}
\text{コストの後の利益} \approx \alpha A - c\,A^{3/2}
\end{equation}
で、利益は $A^* = (2\alpha / 3c)^2$ で最大になり、$A = (\alpha/c)^2$ でゼロになる。予測の減衰が速い（回転率が高い）戦略、小型株を使う戦略ほど、$c$ が大きく容量は小さい。ヘッジファンドが、成績の良い戦略で新しい資金を断ること（ファンドを閉じる）があるのはこのためである。

\section{空売り}

\begin{itemize}
\item \textbf{ロケート}：空売りの前に、借りられる株を証券会社に確保してもらう。
\item \textbf{貸株料}：大半の銘柄では年 0.3\% 程度以下だが、借りたい人の多い銘柄（小型株、悪材料の銘柄）は年数十\% になることもある。空売りで儲かりそうな銘柄ほど借りるのが高い。
\item \textbf{リコールと踏み上げ}：貸し手が株を返せと言えば、空売りを閉じなければならない。空売りの多い銘柄が急に上がると、損失を限るための買い戻しが殺到し、さらに上がる（ショート・スクイーズ）。2021 年 1 月の GameStop が代表例である。
\item \textbf{規制}：危機の時に、空売りが一時的に禁止されることがある（2008 年の金融株など）。
\end{itemize}

\begin{projbox}
LLM-BLM では、片道 0.3\% のコストを、BLM と比べる相手の両方から引いた。小型株の S1 の運用では、上位 20\% を浮動株時価総額で加重し、毎月入れ替える。平方根則で考えると、小さな資金なら 0.3\% は妥当な水準だが、大きな資金では明らかに足りない。逆に、実際に売買する個人や研究の規模では、容量は問題にならない。研究の結論が「この規模の資金で」のものであることを明記する必要がある。
\end{projbox}
